% Table 1: compact dictionary between scaling-law objects and production/IO econometrics (v2, <= 8 rows).
% Writer: framework section writer (v2), 2026-09-24. Replaces the 18-row v1 table (sections/v1/framework.tex).
% Corrections after referee round 1 (R1 comment 8; R2 Major 9a, 9d; R3 M1, minor 10):
%   - on-path collinearity is Marschak-Andrews / Bond-Soderbom collinearity under common prices, not ACF functional
%     dependence (R1 8a); IsoFLOP sweeps are assigned (designed) inputs, not price shocks (R1 8d; R2 9d);
%   - no GNR label for the flagship normalization; no TFPR row for bounded benchmark links (R1 8e, 8h);
%   - Nerlove (1963) is cited for the returns-to-scale (cost-elasticity) measure only, not for mismeasurement;
%     scale-correlated input mismeasurement is attributed to Basu-Fernald (R1 8g);
%   - "Approach 1" renamed: the frontier through IsoFLOP minima is the "IsoFLOP-minima frontier" (R2 9a);
%   - the wedge is Raval's ratio statistic with FLOP-accounting weights (R1 3a).
\begin{table}[tp]
\centering
\caption{Scaling Laws in the Language of Production Economics}
\label{tab:dictionary}
\footnotesize
\setlength{\tabcolsep}{3pt}
\begin{tabular}{@{}>{\raggedright\arraybackslash}p{0.22\textwidth}>{\raggedright\arraybackslash}p{0.22\textwidth}>{\raggedright\arraybackslash}p{0.325\textwidth}>{\raggedright\arraybackslash}p{0.185\textwidth}@{}}
\toprule
Scaling-law object & Production or IO counterpart & Use in this paper & References \\
\midrule
Iso-loss curve; IsoFLOP profile ($6ND=C$) & Isoquant; output along an isocost with unit log-cost weights & The optimum equates output elasticities without price data; interior optima need $\sigma<1$ (Lemma~\ref{lemma:sigma}) & \citealt{shephard1953cost} \\[3pt]
Chinchilla law and its $\kappa$ family & Quasi-homothetic technology; outer exponent $\kappa$ sets returns to scale & All curvature lies across the expansion path; $\kappa=1$ is testable (Section~\ref{sec:tech}) & \citealt{hoffmann2022training}; \citealt{kaplan2020scaling} \\[3pt]
Compute-optimal $N^*(C)$, $D^*(C)$; IsoFLOP-minima frontier & Conditional factor demands; inverse cost function & Duality links $a$, $\gamma$ and $\sigma^*$; the elasticity of cost with respect to reducible loss is $-1/\gamma$ (Proposition~\ref{prop:duality}) & \citealt{uzawa1964duality}; \citealt{nerlove1963returns} \\[3pt]
Compute-optimal ladders; fixed tokens-per-parameter designs & Collinear inputs under common prices & Identify the path and the frontier, not $\sigma^*$; information about curvature is fourth order in allocation errors (Proposition~\ref{prop:ident}) & \citealt{marschak1944random}; \citealt{bond2005adjustment} \\[3pt]
IsoFLOP sweeps; factorial designs & Assigned variation in the input mix; optimal design & Curvature at the IsoFLOP minimum identifies $\sigma^*$ without a functional form (Proposition~\ref{prop:modelfree}) & \citealt{box1959design}; \citealt{kiefer1959optimum} \\[3pt]
Over-training, $M>M^*(C)$ & Input-specific wedge: ratio of output to cost elasticities & Reveals the value of compactness and its expenditure share $s$ (Proposition~\ref{prop:wedge}; Section~\ref{sec:wedge}) & \citealt{deloecker2012markups}; \citealt{raval2023testing} \\[3pt]
Lab-specific recipes, beliefs and output concepts & Factor-augmenting heterogeneity; output-concept mismatch & Contaminate the wedge: $\hat w=w\,w_L/w_{L'}$ (Proposition~\ref{prop:wedge}) & \citealt{demirer2020production}; \citealt{bond2020unpleasant} \\[3pt]
Embedding counts; omitted FLOPs; untuned hyperparameters & Scale-correlated input mismeasurement; flexible inputs not re-optimized & Shift the allocation exponent $a$ and $M^*$ (Section~\ref{sec:tech}) & \citealt{basu1997returns}; \citealt{porian2024resolving} \\
\bottomrule
\end{tabular}
\par\smallskip
\begin{tablenotes}
The first three rows are exact once training compute is $C=6ND$; the others are analogies in which the econometric
problem is the same but its source can differ. $M=D/N$ is tokens per parameter and $M^*(C)$ its compute-optimal value;
$w=\varepsilon_N/\varepsilon_D$ is the wedge and $s=(w-1)/w$; $a$ is the allocation exponent, $\gamma$ the frontier
elasticity, $\sigma^*$ the elasticity of substitution on the expansion path and $\kappa$ the outer exponent; $L$ and
$L'$ are the technologies used by the econometrician and by the developer.
\end{tablenotes}
\begin{tablenotes}[Source]
Authors' compilation.
\end{tablenotes}
\end{table}
